% ================= Phase 3 新增內容 =================
% 4.4～4.8 接在 Section 4 之後；Section 5 為部分功能的實作與展示。

\newcommand{\shot}[3]{% {檔名}{寬度}{說明}
\begin{figure}[htbp]\centering
\IfFileExists{img/#1}{\includegraphics[width=#2\linewidth]{img/#1}}{\fbox{\parbox{0.8\linewidth}{\centering\vspace{2em}（截圖：#1）\vspace{2em}}}}
\caption{#3}\end{figure}}

\needspace{10\baselineskip}
\subsection{預期的資料庫操作、頻率與資料量 (Expectation of the Possible DB Operations, Frequencies and Data Volumes)}

我們以營運一年後的規模估計每日的資料庫操作。估計的前提為：每日約 3,000 位訪客瀏覽網站，其中約 1,500 位顧客登入；每日成立約 300 筆訂單，每筆平均 3 項商品；約 30 位工作人員在後台處理訂單、商品與折扣。

表中「預估使用頻率」為每日透過 SQL 執行該操作的次數；「表格資料量」為該資料表一年後的預估 tuple 數（估算依據見 4.5 節）；「系統負擔」以最壞情況計算，假設每次查詢、修改或刪除都要掃過整張表，即「頻率 × 表格資料量」；新增只寫入一筆，與表格大小無關，因此資料量記為 don't care，負擔即為頻率。

\begingroup\small
\begin{xltabular}{\linewidth}{|>{\raggedright\arraybackslash}p{2.9cm}|>{\raggedright\arraybackslash}X|>{\raggedleft\arraybackslash}p{1.5cm}|>{\raggedleft\arraybackslash}p{1.6cm}|>{\raggedleft\arraybackslash}p{3.2cm}|}
\caption{預期的資料庫操作、頻率與系統負擔}\\
\hline\rowcolor{headgray}\textbf{Table} & \textbf{可能操作} & \textbf{預估使用頻率 (per day)} & \textbf{表格資料量 (tuples)} & \textbf{系統負擔 (worst case，次／day)}\\\hline\endfirsthead
\hline\rowcolor{headgray}\textbf{Table} & \textbf{可能操作} & \textbf{預估使用頻率 (per day)} & \textbf{表格資料量 (tuples)} & \textbf{系統負擔 (worst case，次／day)}\\\hline\endhead
\at{users} & 登入時以 loginid 驗證帳號 & 2,250 & 20,033 & 45,074,250 Query\\\hline
\at{users} & 註冊新顧客 & 30 & don't care & 30 Insert\\\hline
\at{users} & 平台管理員修改資料、停用或啟用帳號 & 20 & 20,033 & 400,660 Update\\\hline
\at{user_role} & 登入時讀取角色 & 2,250 & 20,040 & 45,090,000 Query\\\hline
\at{user_role} & trigger 檢查顧客或工作人員角色 & 4,800 & 20,040 & 96,192,000 Query\\\hline
\at{user_role} & 平台管理員指派或移除角色 & 2 & don't care & 2 Insert\\\hline
\at{store} & 商品頁與商店頁顯示商店資訊 & 6,000 & 10 & 60,000 Query\\\hline
\at{store_staff} & 工作人員操作時查詢隸屬的商店 & 1,000 & 30 & 30,000 Query\\\hline
\at{store_staff} & 平台管理員指定工作人員的商店 & 1 & 30 & 30 Update\\\hline
\at{product} & 關鍵字搜尋、瀏覽與商品詳細資料 & 45,000 & 2,000 & 90,000,000 Query\\\hline
\at{product} & 結帳扣庫存、取消或退貨回補 & 920 & 2,000 & 1,840,000 Update\\\hline
\at{product} & 工作人員新增或修改商品 & 50 & 2,000 & 100,000 Update\\\hline
\at{cart_item} & 加入購物車、修改數量 & 4,500 & don't care & 4,500 Insert\\\hline
\at{cart_item} & 顯示購物車 & 6,000 & 9,000 & 54,000,000 Query\\\hline
\at{cart_item} & 移除商品、結帳後清空 & 1,200 & 9,000 & 10,800,000 Delete\\\hline
\at{orders} & 確認付款後建立訂單 & 300 & don't care & 300 Insert\\\hline
\at{orders} & 顧客查詢歷史訂單、工作人員查詢訂單 & 1,950 & 109,500 & 213,525,000 Query\\\hline
\at{orders} & 更新訂單狀態、取消訂單 & 920 & 109,500 & 100,740,000 Update\\\hline
\at{orders} & 日、週、月銷售報表 & 150 & 109,500 & 16,425,000 Query\\\hline
\at{order_item} & 建立訂單明細 & 900 & don't care & 900 Insert\\\hline
\at{order_item} & 查詢訂單明細、報表統計 & 2,100 & 328,500 & 689,850,000 Query\\\hline
\at{order_status_history} & trigger 自動記錄狀態變更 & 1,220 & don't care & 1,220 Insert\\\hline
\at{order_status_history} & 顧客查看狀態歷程 & 500 & 416,100 & 208,050,000 Query\\\hline
\at{discount} & 購物車預估與結帳時查詢有效折扣 & 6,300 & 200 & 1,260,000 Query\\\hline
\at{discount} & 工作人員新增、修改或停用折扣 & 10 & 200 & 2,000 Update\\\hline
\at{product_discount} & 商品頁顯示特殊活動折扣 & 45,000 & 1,000 & 45,000,000 Query\\\hline
\at{product_discount} & 設定特殊活動適用商品 & 5 & don't care & 5 Insert\\\hline
\at{order_discount} & 記錄訂單套用的折扣 & 450 & don't care & 450 Insert\\\hline
\at{order_discount} & 查詢訂單金額明細 & 750 & 164,250 & 123,187,500 Query\\\hline
\at{return_request} & 申請退貨 & 9 & don't care & 9 Insert\\\hline
\at{return_request} & 查詢退貨申請 & 100 & 3,285 & 328,500 Query\\\hline
\at{return_request} & 審核退貨申請 & 9 & 3,285 & 29,565 Update\\\hline
\at{return_item} & 新增退貨明細、核准時寫入退款金額 & 28 & 4,928 & 137,984 Update\\\hline
\at{store_review} & 撰寫商店評價 & 60 & don't care & 60 Insert\\\hline
\at{store_review} & 商店頁顯示平均評分與評論 & 6,000 & 28,470 & 170,820,000 Query\\\hline
\at{system_setting} & 計算金額時讀取稅率與基本運費 & 6,300 & 2 & 12,600 Query\\\hline

\end{xltabular}
\endgroup

最壞情況下負擔最大的是商品搜尋與商品頁（\at{product}、\at{product_discount}，每日約 9 千萬次 tuple 讀取）、登入與角色檢查（\at{users}、\at{user_role}），以及訂單相關查詢。實際上這些查詢都依主鍵、\at{loginid} 的 UNIQUE 索引或外鍵查找，不需要掃過整張表；Phase 4 會再針對商品關鍵字搜尋與訂單查詢加上索引並調校。

\needspace{10\baselineskip}
\subsection{資料量估計 (Estimated Data Volume)}

每張資料表的大小以「schema 長度 × 預估 tuple 數」估計。schema 長度為一筆資料各欄位平均長度的總和：\at{INT}／\at{SERIAL} 4 bytes、\at{SMALLINT} 2 bytes、\at{NUMERIC} 約 8 bytes、\at{BOOLEAN} 1 byte、\at{DATE} 4 bytes、\at{TIMESTAMP} 8 bytes；\at{VARCHAR} 與 \at{TEXT} 以平均實際長度估計（中文以 UTF-8 每字 3 bytes 計），另加 1 byte 長度標頭。tuple 數為營運一年後的預估筆數。估計不含 PostgreSQL 每筆約 24 bytes 的列標頭與索引空間，實際占用約為估計值的 1.5 至 2 倍。

\begingroup\small
\begin{xltabular}{\linewidth}{|>{\raggedright\arraybackslash}p{3.3cm}|>{\raggedleft\arraybackslash}p{1.9cm}|>{\raggedleft\arraybackslash}p{1.8cm}|>{\raggedleft\arraybackslash}p{1.9cm}|>{\raggedright\arraybackslash}X|}
\caption{各資料表的預估大小}\\
\hline\rowcolor{headgray}\textbf{資料表} & \textbf{schema 長度 (bytes)} & \textbf{tuple 數} & \textbf{預估大小} & \textbf{tuple 數估算依據}\\\hline\endfirsthead
\hline\rowcolor{headgray}\textbf{資料表} & \textbf{schema 長度 (bytes)} & \textbf{tuple 數} & \textbf{預估大小} & \textbf{tuple 數估算依據}\\\hline\endhead
\at{users} & 195 & 20,033 & 3.73 MB & 顧客 20,000、工作人員 30、平台管理員 3\\\hline
\at{user_role} & 11 & 20,040 & 215.3 KB & 每位使用者約 1 個角色，少數兼任\\\hline
\at{store_staff} & 8 & 30 & 0.2 KB & 工作人員約 30 位，大多隸屬一間商店\\\hline
\at{store} & 121 & 10 & 1.2 KB & 10 間門市\\\hline
\at{product} & 446 & 2,000 & 871.1 KB & 每店約 200 項商品\\\hline
\at{cart_item} & 20 & 9,000 & 175.8 KB & 3,000 位顧客購物車內平均 3 項\\\hline
\at{orders} & 155 & 109,500 & 16.19 MB & 每日 300 筆，保存一年\\\hline
\at{order_item} & 67 & 328,500 & 20.99 MB & 每筆訂單平均 3 項\\\hline
\at{order_status_history} & 27 & 416,100 & 10.71 MB & 每筆訂單平均 3.8 次狀態\\\hline
\at{discount} & 114 & 200 & 22.3 KB & 一年約 200 個折扣政策\\\hline
\at{product_discount} & 37 & 1,000 & 36.1 KB & 特殊活動平均套用 5 項商品\\\hline
\at{order_discount} & 23 & 164,250 & 3.60 MB & 每筆訂單平均套用 1.5 個折扣\\\hline
\at{return_request} & 137 & 3,285 & 439.5 KB & 約 3\% 的訂單申請退貨\\\hline
\at{return_item} & 22 & 4,928 & 105.9 KB & 每筆申請平均退 1.5 項\\\hline
\at{store_review} & 109 & 28,470 & 2.96 MB & 約 20\% 訂單評價、平均 1.3 間店\\\hline
\at{system_setting} & 52 & 2 & 0.1 KB & 稅率、基本運費\\\hline
\multicolumn{3}{|r|}{\textbf{合計}} & \textbf{60.00 MB} & 一年資料量\\\hline

\end{xltabular}
\endgroup

一年的資料量約 60 MB，最大的是 \at{order_item}、\at{orders} 與 \at{order_status_history}，三者都會隨訂單數線性成長；商品、商店、折扣等主檔資料量很小。

\clearpage
\subsection{建立資料庫的 SQL 語法 (SQL Statements Used to Construct the Schema)}

資料庫建立在 PostgreSQL 16 上，SQL 分成四個檔案，依序執行：

\begin{enumerate}[itemsep=0.2em]
  \item \at{01_schema.sql}：建立 16 張資料表與主鍵、外鍵、UNIQUE、CHECK、EXCLUDE 限制，內容與 4.1、4.2 節一致。
  \item \at{02_triggers.sql}：建立 4.3 節表中的 trigger 與輔助函式。
  \item \at{03_seed.sql}：匯入系統運作需要的初始資料（4.7 節）。
  \item \at{04_demo_data.sql}：匯入示範用的交易資料（4.7 節）。
\end{enumerate}

\textbf{與 Phase 2 的差異。}\at{return_request} 新增 \at{refund_total} 欄位。Phase 2 的 ER 圖把退款總額畫成衍生屬性；實作時退貨核准後退款金額就固定，不會再變動，顧客與工作人員查詢退貨結果時也常用到，所以在核准當下依 BR-26 算好後存起來，避免每次查詢都要加總 \at{return_item}。

\subsubsection{資料表與限制（01\_schema.sql）}
\VerbatimInput[fontsize=\scriptsize,numbers=left,numbersep=4pt]{sql/01_schema.sql}

\subsubsection{Trigger（02\_triggers.sql）}
\VerbatimInput[fontsize=\scriptsize,numbers=left,numbersep=4pt]{sql/02_triggers.sql}

\clearpage
\subsection{匯入資料的 SQL 語法 (SQL Statements Used to Insert the Data)}

\at{03_seed.sql} 匯入系統運作需要的初始資料（DBS-F-002），包含系統參數、平台管理員與工作人員帳號、示範顧客、商店、商品與折扣政策。密碼以 bcrypt 雜湊後儲存。

\VerbatimInput[fontsize=\scriptsize,numbers=left,numbersep=4pt]{sql/03_seed.sql}

\at{04_demo_data.sql} 匯入示範用的交易資料：13 筆不同狀態的訂單與其明細、套用的折扣、退貨申請與商店評價，以及購物車內容。訂單金額依 BR-09～BR-12 計算；狀態歷程由 trigger 自動寫入；所有資料都通過 4.3 節的 trigger 檢查。

\VerbatimInput[fontsize=\scriptsize,numbers=left,numbersep=4pt]{sql/04_demo_data.sql}

\clearpage
\subsection{資料表在 PostgreSQL 中的實作 (Implementation of Tables in the Target DBMS)}

我們以 Docker 啟動 PostgreSQL 16（\at{postgres:16-alpine}），依序執行 \at{sql/} 中的 SQL 檔建立資料庫並匯入資料。以下截圖為在 \at{psql} 中查詢的結果。

\newcommand{\dbshot}[3]{% {檔名}{寬度}{說明}
\begin{figure}[htbp]\centering\includegraphics[width=#2\linewidth]{img/#1}\caption{#3}\end{figure}}

\textbf{資料表清單。}資料庫中共有 4.1 節的 16 張資料表（圖 \ref{fig:db1}）。執行 \at{03_seed.sql} 與 \at{04_demo_data.sql} 後，各表的資料筆數如圖 \ref{fig:db2}：示範資料包含 3 間商店、15 項商品、9 位使用者（含 3 位隸屬商店的工作人員），以及 13 筆涵蓋 received、processing、shipping、closed、cancelled、returned 各種狀態的訂單。

\dbshot{db1_tables.png}{0.85}{\label{fig:db1}資料庫中的 16 張資料表（\texttt{\textbackslash dt+}）}
\dbshot{db2_counts.png}{0.5}{\label{fig:db2}各資料表的資料筆數}

\textbf{資料表結構。}以 \texttt{\textbackslash d} 查看資料表定義，可以看到欄位型態、預設值、主鍵、CHECK、外鍵、被哪些表參照，以及掛在表上的 trigger。圖 \ref{fig:db3} 為 \at{orders}；圖 \ref{fig:db4} 為 \at{product_discount}，其中可以看到實作 BR-13 的 EXCLUDE 限制與參照 \at{discount} 的複合外鍵。

\dbshot{db3_orders.png}{0.95}{\label{fig:db3}\at{orders} 的結構、限制與 trigger}
\dbshot{db4_product_discount.png}{0.95}{\label{fig:db4}\at{product_discount} 的 EXCLUDE 限制與複合外鍵}

\textbf{資料內容。}圖 \ref{fig:db5} 為示範訂單的金額欄位，皆依 BR-09～BR-12 計算；圖 \ref{fig:db6} 為訂單明細，以及由 trigger 自動寫入的訂單狀態歷程；圖 \ref{fig:db7} 為退貨申請與退貨明細，包含已核准、已拒絕與待審核三種狀態。

\dbshot{db5_orders_data.png}{0.9}{\label{fig:db5}\at{orders} 的資料}
\dbshot{db6_items.png}{0.85}{\label{fig:db6}\at{order_item} 與 \at{order_status_history} 的資料}
\dbshot{db7_returns.png}{0.9}{\label{fig:db7}\at{return_request} 與 \at{return_item} 的資料}

\textbf{完整性限制的驗證。}直接以 SQL 寫入違反商業規則的資料，資料庫都會拒絕（圖 \ref{fig:db8}）：訂單狀態從 received 直接跳到 closed 被 trigger 拒絕（BR-19）；把顧客 alice 設為商店工作人員被 trigger \at{store_staff_check} 拒絕（BR-21）；同一商品設定期間重疊的特殊活動折扣被 EXCLUDE 拒絕（BR-13）；購物車數量為 0 被 CHECK 拒絕（BR-05）；同一筆訂單第二次申請退貨被 UNIQUE 拒絕（BR-25）。

\dbshot{db8_reject.png}{0.95}{\label{fig:db8}違反商業規則的資料被資料庫拒絕}


% ================= Section 5 =================
\clearpage
\section{資料庫系統實作與展示 (Implementation and Demonstration)}

\subsection{系統安裝}
系統由三個 Docker 容器組成：PostgreSQL 16 資料庫、FastAPI 後端（Python，JWT 登入、SQLAlchemy 存取資料庫）與 React + Vite 前端。資料庫第一次啟動時，會依檔名順序自動執行 \at{db/init/} 中的四個 SQL 檔（即 Appendix 的 SQL）。

\textbf{安裝需求。}安裝 Docker Desktop 即可，不需要另外安裝 PostgreSQL、Python 或 Node.js。

\textbf{啟動步驟。}在專案根目錄執行：
\begin{Verbatim}[fontsize=\small,xleftmargin=2em]
./start.sh            # 等同 docker compose up --build
./start.sh reset      # 清空資料庫，重新執行 db/init/*.sql
./start.sh down       # 停止所有容器
\end{Verbatim}

\begingroup\small
\begin{xltabular}{\linewidth}{|>{\raggedright\arraybackslash}p{4cm}|>{\raggedright\arraybackslash}X|}
\caption{系統服務與網址}\\
\hline\rowcolor{headgray}\textbf{服務} & \textbf{網址}\\\hline\endhead
網站（三種介面） & \texttt{http://localhost:5173}\\\hline
後端 API 文件（Swagger） & \texttt{http://localhost:8001/docs}\\\hline
PostgreSQL & \texttt{localhost:5433}，帳號、密碼、資料庫名稱皆為 \texttt{fuwa}\\\hline
\end{xltabular}
\endgroup

\textbf{示範帳號。}\at{03_seed.sql} 建立了以下帳號，登入後依角色進入不同的介面：

\begingroup\small
\begin{xltabular}{\linewidth}{|>{\raggedright\arraybackslash}p{2cm}|>{\raggedright\arraybackslash}p{2.4cm}|>{\raggedright\arraybackslash}p{3.4cm}|>{\raggedright\arraybackslash}X|}
\caption{示範帳號}\\
\hline\rowcolor{headgray}\textbf{帳號} & \textbf{密碼} & \textbf{角色} & \textbf{進入的介面}\\\hline\endhead
\texttt{admin} & \texttt{admin123} & 平台管理員 & 平台管理介面（\texttt{/admin}）\\\hline
\texttt{staff} & \texttt{staff123} & 工作人員 & 工作人員管理介面（\texttt{/staff}）\\\hline
\texttt{manager} & \texttt{staff123} & 工作人員、平台管理員 & 兩個介面都能進入，頁首可切換\\\hline
\texttt{alice}、\texttt{bob} & \texttt{customer123} & 顧客 & 顧客購物介面\\\hline
\texttt{carol} & \texttt{customer123} & 顧客（已停用） & 無法登入，用來示範停用帳號\\\hline
\end{xltabular}
\endgroup


\subsection{系統展示}
\newcommand{\guishot}[2]{% {檔名}{說明}
\begin{figure}[htbp]\centering\includegraphics[width=0.92\linewidth]{img/#1}\caption{#2}\end{figure}}

本節展示系統的核心流程：顧客瀏覽商品、加入購物車並查看訂單，以及工作人員處理訂單。以下截圖為系統以 4.7 節的示範資料啟動後實際操作的畫面；退貨審核、折扣管理、銷售報表與平台管理等其餘功能將在 Phase 4 展示。

\subsubsection*{顧客購物介面}

未登入的訪客也可以瀏覽商品。首頁（圖 \ref{fig:gui_home}）上方列出今天有效的折扣政策，可以依關鍵字、商店與類別搜尋商品；點選商品進入商品頁（圖 \ref{fig:gui_product}）。

\guishot{gui_home.png}{\label{fig:gui_home}首頁：今日優惠、商品搜尋與篩選}
\guishot{gui_product.png}{\label{fig:gui_product}商品頁}

登入頁（圖 \ref{fig:gui_login}）以帳號與密碼登入，已停用的帳號無法登入。顧客登入後可以把商品加入購物車並調整數量（圖 \ref{fig:gui_cart}），結帳後在「我的訂單」查看所有訂單與消費總額（圖 \ref{fig:gui_my_orders}）；訂單詳情顯示訂單進度與狀態歷程（圖 \ref{fig:gui_order_detail}），狀態歷程即 \at{order_status_history} 中由 trigger 寫入的資料。

\guishot{gui_login.png}{\label{fig:gui_login}登入頁}
\guishot{gui_cart.png}{\label{fig:gui_cart}購物車（顧客 alice）}
\guishot{gui_my_orders.png}{\label{fig:gui_my_orders}我的訂單}
\guishot{gui_order_detail.png}{\label{fig:gui_order_detail}訂單詳情與狀態歷程（訂單 \#5）}

\subsubsection*{工作人員管理介面}

工作人員在訂單管理（圖 \ref{fig:gui_staff_orders}）依狀態、日期與關鍵字查詢訂單，並在訂單詳情（圖 \ref{fig:gui_staff_order_detail}）依 BR-19 的順序推進狀態。圖中訂單 \#10 有待審核的退貨申請，因此依 BR-25 不能推進狀態，畫面上方顯示提示。

\guishot{gui_staff_orders.png}{\label{fig:gui_staff_orders}訂單管理}
\guishot{gui_staff_order_detail.png}{\label{fig:gui_staff_order_detail}訂單詳情：有待審核退貨時不能推進狀態（BR-25）}


